% Exercice: Probleme 11
% Sujet: Algèbre
% Généré automatiquement

\begin{exercice}[Algèbre]
\label{ex:59}

Problème\\
1^\{1\} points\\
Partle A\\
3,5 polnts\\
1^\{0\} 0^\{00\}\\
x?\\
définie par f(x)\\
est la fonction numérique dune variable réelle\\
7^\{3\})\\
On prend Icm\\
est sa courbe représentative dans le repère orthogonal (0,\\
pour 5^\{0\} sur l'axe des abscisses et Icm\\
0^\{00\} sur l'axe des ordonnées\\
pour\\
Etudier les variations de\\
sur [-1^\{00\}\\
1^\{00\}] et dresser son tableau de variations\_\\
2^\{0\} Construire la courbe Cf pour les abscisses appartenant\\
1^\{00\}]\\
Vintervalle [-1^\{00\}\\
définie par g(x)\\
f(x) + 1\\
3^\{0\} Soit\\
la fonction numérique de la variable réelle\\
$(0,7,3)$\\
Tracer la courbe Cg de g dans le repère\\
(On indiquera le programme de construction de C, à partir de Cf)\\

\end{exercice}

% --- Fin Probleme 11 ---
